% GENERATED FILE. Edit canonical lessons, then run npm run generate:books.
% canonical-source-hash: fnv1a64:b049a7fd32bc2e5a
% canonical-lessons: MR-C428-udhalne, MR-C428-sillak, MR-C428-guntavnuk, MR-C428-aipat, MR-C428-kiphayatsir

\chapter{Spending and Saving}
\label{ch:mr-spending-and-saving}
\hlchaptermodality{428}

\begin{chapteropening}
\textbf{By the end of this chapter:} \emph{I can talk about squandering, a balance, an investment, what I can afford and a thrifty choice.}

Complete the last lesson of chapter 428: I can talk about squandering, a balance, an investment, what I can afford and a thrifty choice.
\end{chapteropening}

\section[\texorpdfstring{\mr{उधळणे} (udhaḷṇe)}{उधळणे (udhaḷṇe)}]{\mr{उधळणे} (udhaḷṇe) — to squander}

\label{lesson:MR-C428-udhalne}



\begin{quote}
Before the new one: say the Marathi for the remainder, then the Marathi for to repay.
\end{quote}

\subsection*{You'll want to know: \mr{उधळणे}}
\textbf{\mr{उधळणे}} — \emph{udhaḷṇe} — \textquotedblleft{}to squander\textquotedblright{}.

\begin{quote}
\mr{त्याने} \mr{सगळा} \mr{पगार} \mr{एका} \mr{आठवड्यात} \mr{उधळला}.
\end{quote}

\begin{quote}
\emph{tyāne sagḷā pagār ekā āṭhvaḍyāt udhaḷlā.}
\end{quote}

\begin{quote}
He squandered his whole salary in one week.
\end{quote}

Past, with -\mr{ने} (he …): \mr{त्याने} \mr{उधळलं}, \emph{tyāne udhaḷlaṁ}.

Also to scatter: \mr{फुलं} \mr{उधळणे}, to shower flowers.

\begin{cousinweb}
\textbf{\mr{उधळणे}} (\emph{udhaḷṇe}) is inherited: CDIAL lists Marathi \emph{udhaḷṇẽ}, \textquotedblleft{}to scatter; to spend lavishly\textquotedblright{}, under Sanskrit \emph{uddhūlayati}, \textquotedblleft{}sprinkles with powder\textquotedblright{} (CDIAL 2027).
\end{cousinweb}

\begin{grammarlens}[title={What you've built}]
One more word for getting around.
\end{grammarlens}

\subsection*{Your turn}
\begin{itemize}
\raggedright
  \item \emph{Say it:} \emph{udhaḷṇe}
  \item \emph{Say it:} \emph{udhaḷṇe}, once more
  \item \emph{Say it:} \emph{tyāne sagḷā pagār ekā āṭhvaḍyāt udhaḷlā.} — He squandered his whole salary in one week.
  \item \emph{Recall:} say the Marathi for the remainder, then the Marathi for to repay, then say \emph{udhaḷṇe} again
\end{itemize}

\begin{tcolorbox}[breakable,colback=teal!4,colframe=teal!35!black,title={Before you move on}]
What does \mr{उधळणे} mean? (\textquotedblleft{}To squander\textquotedblright{}.) Say it once more.
\end{tcolorbox}
\section[\texorpdfstring{\mr{शिल्लक} (śillak)}{शिल्लक (śillak)}]{\mr{शिल्लक} (śillak) — a balance; what is left}

\label{lesson:MR-C428-sillak}



\begin{quote}
Before the new one: say the Marathi for to repay, then the Marathi for to squander.
\end{quote}

\subsection*{You'll want to know: \mr{शिल्लक}}
\textbf{\mr{शिल्लक}} — \emph{śillak} — \textquotedblleft{}a balance; what is left\textquotedblright{}.

\begin{quote}
\mr{महिनाअखेर} \mr{माझ्याकडे} \mr{थोडी} \mr{शिल्लक} \mr{राहते}.
\end{quote}

\begin{quote}
\emph{mahināakher mājhyākaḍe thoḍī śillak rāhte.}
\end{quote}

\begin{quote}
At the end of the month I have a little left over.
\end{quote}

Feminine.

\mr{शिल्लक} \mr{राहणे}, to be left over; \mr{खात्यातली} \mr{शिल्लक}, the account balance.

\begin{cousinweb}
\textbf{\mr{शिल्लक}} (\emph{śillak}) is a Marathi word for what remains of money or stock; its origin is not settled.
\end{cousinweb}

\begin{grammarlens}[title={What you've built}]
One more word for getting around.
\end{grammarlens}

\subsection*{Your turn}
\begin{itemize}
\raggedright
  \item \emph{Say it:} \emph{śillak}
  \item \emph{Say it:} \emph{śillak}, once more
  \item \emph{Say it:} \emph{mahināakher mājhyākaḍe thoḍī śillak rāhte.} — At the end of the month I have a little left over.
  \item \emph{Recall:} say the Marathi for to repay, then the Marathi for to squander, then say \emph{śillak} again
\end{itemize}

\begin{tcolorbox}[breakable,colback=teal!4,colframe=teal!35!black,title={Before you move on}]
What does \mr{शिल्लक} mean? (\textquotedblleft{}A balance; what is left\textquotedblright{}.) Say it once more.
\end{tcolorbox}
\section[\texorpdfstring{\mr{गुंतवणूक} (guntavṇūk)}{गुंतवणूक (guntavṇūk)}]{\mr{गुंतवणूक} (guntavṇūk) — an investment}

\label{lesson:MR-C428-guntavnuk}



\begin{quote}
Before the new one: say the Marathi for to squander, then the Marathi for a balance.
\end{quote}

\subsection*{You'll want to know: \mr{गुंतवणूक}}
\textbf{\mr{गुंतवणूक}} — \emph{guntavṇūk} — \textquotedblleft{}an investment\textquotedblright{}.

\begin{quote}
\mr{घर} \mr{ही} \mr{चांगली} \mr{गुंतवणूक} \mr{आहे}.
\end{quote}

\begin{quote}
\emph{ghar hī cāṅglī guntavṇūk āhe.}
\end{quote}

\begin{quote}
A house is a good investment.
\end{quote}

Feminine.

\mr{गुंतवणूक} \mr{करणे}, to invest; a feminine noun.

\begin{cousinweb}
\textbf{\mr{गुंतवणूक}} (\emph{guntavṇūk}) is built on \textbf{\mr{गुंतवणे}}, \textquotedblleft{}to entangle; to invest\textquotedblright{}, with the noun ending -\mr{णूक}; the origin of the verb is not settled.
\end{cousinweb}

\begin{grammarlens}[title={What you've built}]
One more word for getting around.
\end{grammarlens}

\subsection*{Your turn}
\begin{itemize}
\raggedright
  \item \emph{Say it:} \emph{guntavṇūk}
  \item \emph{Say it:} \emph{guntavṇūk}, once more
  \item \emph{Say it:} \emph{ghar hī cāṅglī guntavṇūk āhe.} — A house is a good investment.
  \item \emph{Recall:} say the Marathi for to squander, then the Marathi for a balance, then say \emph{guntavṇūk} again
\end{itemize}

\begin{tcolorbox}[breakable,colback=teal!4,colframe=teal!35!black,title={Before you move on}]
What does \mr{गुंतवणूक} mean? (\textquotedblleft{}An investment\textquotedblright{}.) Say it once more.
\end{tcolorbox}
\section[\texorpdfstring{\mr{ऐपत} (aipat)}{ऐपत (aipat)}]{\mr{ऐपत} (aipat) — what one can afford; means}

\label{lesson:MR-C428-aipat}



\begin{quote}
Before the new one: say the Marathi for a balance, then the Marathi for an investment.
\end{quote}

\subsection*{You'll want to know: \mr{ऐपत}}
\textbf{\mr{ऐपत}} — \emph{aipat} — \textquotedblleft{}what one can afford; means\textquotedblright{}.

\begin{quote}
\mr{घर} \mr{घेण्याइतकी} \mr{माझी} \mr{ऐपत} \mr{नाही}.
\end{quote}

\begin{quote}
\emph{ghar gheṇyāitkī mājhī aipat nāhī.}
\end{quote}

\begin{quote}
I cannot afford to buy a house.
\end{quote}

Feminine.

Usually with the owner in -\mr{ची}: \mr{माझी} \mr{ऐपत}, my means.

\begin{cousinweb}
\textbf{\mr{ऐपत}} (\emph{aipat}) is a Marathi word for how much money one has to spend; its origin is not settled.
\end{cousinweb}

\begin{grammarlens}[title={What you've built}]
One more word for getting around.
\end{grammarlens}

\subsection*{Your turn}
\begin{itemize}
\raggedright
  \item \emph{Say it:} \emph{aipat}
  \item \emph{Say it:} \emph{aipat}, once more
  \item \emph{Say it:} \emph{ghar gheṇyāitkī mājhī aipat nāhī.} — I cannot afford to buy a house.
  \item \emph{Recall:} say the Marathi for a balance, then the Marathi for an investment, then say \emph{aipat} again
\end{itemize}

\begin{tcolorbox}[breakable,colback=teal!4,colframe=teal!35!black,title={Before you move on}]
What does \mr{ऐपत} mean? (\textquotedblleft{}What one can afford; means\textquotedblright{}.) Say it once more.
\end{tcolorbox}
\section[\texorpdfstring{\mr{किफायतशीर} (kiphāyatśīr)}{किफायतशीर (kiphāyatśīr)}]{\mr{किफायतशीर} (kiphāyatśīr) — economical; good value}

\label{lesson:MR-C428-kiphayatsir}



\begin{quote}
Before the new one: say the Marathi for an investment, then the Marathi for what one can afford.
\end{quote}

\subsection*{You'll want to know: \mr{किफायतशीर}}
\textbf{\mr{किफायतशीर}} — \emph{kiphāyatśīr} — \textquotedblleft{}economical; good value\textquotedblright{}.

\begin{quote}
\mr{बसचा} \mr{प्रवास} \mr{किफायतशीर} \mr{असतो}.
\end{quote}

\begin{quote}
\emph{bascā pravās kiphāyatśīr asto.}
\end{quote}

\begin{quote}
Travelling by bus is economical.
\end{quote}

It does not change: \mr{किफायतशीर} \mr{दर}, \mr{किफायतशीर} \mr{गाडी}.

\begin{cousinweb}
\textbf{\mr{किफायतशीर}} (\emph{kiphāyatśīr}) is a loan built on Arabic \emph{kifāyat}, \textquotedblleft{}sufficiency, thrift\textquotedblright{}, with the Persian adjective ending \emph{-śīr}; the exact path into Marathi is not settled.
\end{cousinweb}

\begin{grammarlens}[title={What you've built}]
One more word for getting around.
\end{grammarlens}

\subsection*{Your turn}
\begin{itemize}
\raggedright
  \item \emph{Say it:} \emph{kiphāyatśīr}
  \item \emph{Say it:} \emph{kiphāyatśīr}, once more
  \item \emph{Say it:} \emph{bascā pravās kiphāyatśīr asto.} — Travelling by bus is economical.
  \item \emph{Recall:} say the Marathi for an investment, then the Marathi for what one can afford, then say \emph{kiphāyatśīr} again
\end{itemize}

\begin{tcolorbox}[breakable,colback=teal!4,colframe=teal!35!black,title={Before you move on}]
What does \mr{किफायतशीर} mean? (\textquotedblleft{}Economical; good value\textquotedblright{}.) Say it once more.
\end{tcolorbox}
